\chapter{What's New In Bijou v4.1.0}\label{record-275}

\noindent{\footnotesize\color{BookMuted}\path{docs/releases/4.1.0/whats-new.md}\par}\smallskip

This is the long-form release overview for Bijou \texttt{v\allowbreak{}4\allowbreak{}.\allowbreak{}1\allowbreak{}.\allowbreak{}0}.

It is intentionally different from \href{https://github.com/flyingrobots/bijou/blob/cbd70612140e7a450d764b3c7be1c6af878b99bc/docs/CHANGELOG.md}{\texttt{C\allowbreak{}H\allowbreak{}A\allowbreak{}N\allowbreak{}G\allowbreak{}E\allowbreak{}L\allowbreak{}O\allowbreak{}G\allowbreak{}.\allowbreak{}m\allowbreak{}d}}: the changelog is the ledger, while this document explains the shape and meaning of the release.

\section{Release Story}\label{record-275-release-story}

The short version is:

Bijou has moved from a surface-first terminal toolkit toward a more coherent application platform.

The main themes are:

\begin{itemize}
\tightlist
\item
  the runtime now has a real engine model instead of shell folklore
\item
  the framed shell is calmer, richer, and easier to inspect
\item
  DOGFOOD is now a serious proving surface rather than a partial demo
\item
  localization has a real substrate and tooling story
\item
  design language and delivery doctrine are now explicit repo truth
\end{itemize}

If \texttt{v\allowbreak{}4\allowbreak{}.\allowbreak{}0\allowbreak{}.\allowbreak{}0} was the ``pure v4 surface'' release, \texttt{v\allowbreak{}4\allowbreak{}.\allowbreak{}1\allowbreak{}.\allowbreak{}0} is the ``runtime, shell, docs, and tooling grow up together'' release.

\section{Runtime Engine: From Seams To Model}\label{record-275-runtime-engine-from-seams-to-model}

The biggest architectural change since \texttt{v\allowbreak{}4\allowbreak{}.\allowbreak{}0\allowbreak{}.\allowbreak{}0} is the new runtime engine direction captured in:

\begin{itemize}
\tightlist
\item
  \href{https://github.com/flyingrobots/bijou/blob/cbd70612140e7a450d764b3c7be1c6af878b99bc/docs/design/RE-001-define-runtime-engine-architecture.md}{RE-001}
\item
  \href{https://github.com/flyingrobots/bijou/blob/cbd70612140e7a450d764b3c7be1c6af878b99bc/docs/design/RE-002-promote-first-class-state-machine-and-view-stack.md}{RE-002}
\item
  \href{https://github.com/flyingrobots/bijou/blob/cbd70612140e7a450d764b3c7be1c6af878b99bc/docs/design/RE-003-retain-layout-trees-and-layout-invalidation.md}{RE-003}
\item
  \href{https://github.com/flyingrobots/bijou/blob/cbd70612140e7a450d764b3c7be1c6af878b99bc/docs/design/RE-004-route-input-through-layouts-and-layer-bubbling.md}{RE-004}
\item
  \href{https://github.com/flyingrobots/bijou/blob/cbd70612140e7a450d764b3c7be1c6af878b99bc/docs/design/RE-005-buffer-commands-and-effects-separately.md}{RE-005}
\item
  \href{https://github.com/flyingrobots/bijou/blob/cbd70612140e7a450d764b3c7be1c6af878b99bc/docs/design/RE-006-formalize-component-layout-and-interaction-contracts.md}{RE-006}
\end{itemize}

In practical terms, the runtime now exposes a clearer model for:

\begin{itemize}
\tightlist
\item
  state machines versus view stacks
\item
  retained layouts and invalidation causes
\item
  layout-driven input routing
\item
  distinct command and effect buffers
\item
  component-owned interaction contracts
\end{itemize}

This matters because builders can now explain where state lives, where views live, how geometry is retained, and how input is routed without relying on branch order inside the shell. Agents benefit even more: runtime ownership is increasingly inspectable instead of implicit.

The remaining follow-on is \href{https://github.com/flyingrobots/bijou/blob/cbd70612140e7a450d764b3c7be1c6af878b99bc/docs/method/graveyard/legacy-backlog/up-next/RE-007-migrate-framed-shell-onto-runtime-engine-seams.md}{RE-007}, which is about moving more of the framed shell onto those explicit seams and is intentionally deferred until after \texttt{v\allowbreak{}4\allowbreak{}.\allowbreak{}1\allowbreak{}.\allowbreak{}0}.

\section{Framed Shell: More Humane, More Truthful}\label{record-275-framed-shell-more-humane-more-truthful}

The shell-focused work since \texttt{v\allowbreak{}4\allowbreak{}.\allowbreak{}0\allowbreak{}.\allowbreak{}0} substantially improves the framed application experience. The key cycles are:

\begin{itemize}
\tightlist
\item
  \href{https://github.com/flyingrobots/bijou/blob/cbd70612140e7a450d764b3c7be1c6af878b99bc/docs/design/HT-002-layered-focus-and-interaction.md}{HT-002}
\item
  \href{https://github.com/flyingrobots/bijou/blob/cbd70612140e7a450d764b3c7be1c6af878b99bc/docs/design/HT-003-implement-layer-stack-and-input-map-routing.md}{HT-003}
\item
  \href{https://github.com/flyingrobots/bijou/blob/cbd70612140e7a450d764b3c7be1c6af878b99bc/docs/design/HT-004-promote-explicit-layer-objects-and-richer-shell-introspection.md}{HT-004}
\end{itemize}

Notable outcomes:

\begin{itemize}
\tightlist
\item
  framed apps now use a clearer topmost-layer model for routing and dismissal
\item
  help, settings, notifications, quit confirm, and page modals behave more like real shell layers and less like ad hoc branches
\item
  footer and help affordances are more truthful to the actual active layer
\item
  richer shell introspection is exported through helpers like \texttt{d\allowbreak{}e\allowbreak{}s\allowbreak{}c\allowbreak{}r\allowbreak{}i\allowbreak{}b\allowbreak{}e\allowbreak{}F\allowbreak{}r\allowbreak{}a\allowbreak{}m\allowbreak{}e\allowbreak{}L\allowbreak{}a\allowbreak{}y\allowbreak{}e\allowbreak{}r\allowbreak{}S\allowbreak{}t\allowbreak{}a\allowbreak{}c\allowbreak{}k\allowbreak{}(\allowbreak{}.\allowbreak{}.\allowbreak{}.\allowbreak{})}, \texttt{a\allowbreak{}c\allowbreak{}t\allowbreak{}i\allowbreak{}v\allowbreak{}e\allowbreak{}F\allowbreak{}r\allowbreak{}a\allowbreak{}m\allowbreak{}e\allowbreak{}L\allowbreak{}a\allowbreak{}y\allowbreak{}e\allowbreak{}r\allowbreak{}(\allowbreak{}.\allowbreak{}.\allowbreak{}.\allowbreak{})}, and \texttt{u\allowbreak{}n\allowbreak{}d\allowbreak{}e\allowbreak{}r\allowbreak{}l\allowbreak{}y\allowbreak{}i\allowbreak{}n\allowbreak{}g\allowbreak{}F\allowbreak{}r\allowbreak{}a\allowbreak{}m\allowbreak{}e\allowbreak{}L\allowbreak{}a\allowbreak{}y\allowbreak{}e\allowbreak{}r\allowbreak{}(\allowbreak{}.\allowbreak{}.\allowbreak{}.\allowbreak{})}
\end{itemize}

For operators, the result is calmer and safer fullscreen behavior. For builders and agents, the shell is easier to reason about because it advertises more of its actual state.

\section{Framed-App Typing And Command Lifecycles}\label{record-275-framed-app-typing-and-command-lifecycles}

The developer-experience work since \texttt{v\allowbreak{}4\allowbreak{}.\allowbreak{}0\allowbreak{}.\allowbreak{}0} is not just cosmetic. Two important seams got more honest:

\begin{itemize}
\tightlist
\item
  \href{https://github.com/flyingrobots/bijou/blob/cbd70612140e7a450d764b3c7be1c6af878b99bc/docs/design/DX-001-type-framed-app-messages-and-updates-end-to-end.md}{DX-001}
\item
  \href{https://github.com/flyingrobots/bijou/blob/cbd70612140e7a450d764b3c7be1c6af878b99bc/docs/design/DX-002-reconcile-cmd-typing-with-cleanup-and-effect-patterns.md}{DX-002}
\end{itemize}

The framed-shell surface now exports explicit typed helpers such as:

\begin{itemize}
\tightlist
\item
  \texttt{F\allowbreak{}r\allowbreak{}a\allowbreak{}m\allowbreak{}e\allowbreak{}P\allowbreak{}a\allowbreak{}g\allowbreak{}e\allowbreak{}M\allowbreak{}s\allowbreak{}g}
\item
  \texttt{F\allowbreak{}r\allowbreak{}a\allowbreak{}m\allowbreak{}e\allowbreak{}P\allowbreak{}a\allowbreak{}g\allowbreak{}e\allowbreak{}U\allowbreak{}p\allowbreak{}d\allowbreak{}a\allowbreak{}t\allowbreak{}e\allowbreak{}R\allowbreak{}e\allowbreak{}s\allowbreak{}u\allowbreak{}l\allowbreak{}t}
\item
  \texttt{F\allowbreak{}r\allowbreak{}a\allowbreak{}m\allowbreak{}e\allowbreak{}d\allowbreak{}A\allowbreak{}p\allowbreak{}p\allowbreak{}M\allowbreak{}s\allowbreak{}g}
\item
  \texttt{F\allowbreak{}r\allowbreak{}a\allowbreak{}m\allowbreak{}e\allowbreak{}d\allowbreak{}A\allowbreak{}p\allowbreak{}p}
\item
  \texttt{F\allowbreak{}r\allowbreak{}a\allowbreak{}m\allowbreak{}e\allowbreak{}d\allowbreak{}A\allowbreak{}p\allowbreak{}p\allowbreak{}U\allowbreak{}p\allowbreak{}d\allowbreak{}a\allowbreak{}t\allowbreak{}e\allowbreak{}R\allowbreak{}e\allowbreak{}s\allowbreak{}u\allowbreak{}l\allowbreak{}t}
\end{itemize}

And \texttt{C\allowbreak{}m\allowbreak{}d\allowbreak{}<\allowbreak{}M\allowbreak{}>} now reflects the runtime patterns Bijou actually uses. A command may:

\begin{itemize}
\tightlist
\item
  complete synchronously or asynchronously
\item
  emit intermediate messages
\item
  resolve to a final message
\item
  resolve to \texttt{Q\allowbreak{}U\allowbreak{}I\allowbreak{}T}
\item
  return a cleanup handle or cleanup function for long-lived work
\item
  resolve to \texttt{v\allowbreak{}o\allowbreak{}i\allowbreak{}d}
\end{itemize}

That makes the public API line up better with real pulse timers, runtime subscriptions, and sub-app command mapping.

\section{DOGFOOD: From Partial Demo To Real Field Guide}\label{record-275-dogfood-from-partial-demo-to-real-field-guide}

The DOGFOOD work since \texttt{v\allowbreak{}4\allowbreak{}.\allowbreak{}0\allowbreak{}.\allowbreak{}0} is one of the clearest visible changes.

The field guide moved through the coverage ratchet from early partial coverage to complete coverage of all canonical component families, culminating in \href{https://github.com/flyingrobots/bijou/blob/cbd70612140e7a450d764b3c7be1c6af878b99bc/docs/design/DF-019-raise-dogfood-coverage-floor-to-100-percent.md}{DF-019}.

What changed:

\begin{itemize}
\tightlist
\item
  DOGFOOD now covers \texttt{3\allowbreak{}5\allowbreak{}/\allowbreak{}3\allowbreak{}5} canonical component families
\item
  the coverage floor is enforced by repo tooling instead of being a decorative progress bar
\item
  the docs app in \href{https://github.com/flyingrobots/bijou/blob/cbd70612140e7a450d764b3c7be1c6af878b99bc/examples/docs/README.md}{\texttt{e\allowbreak{}x\allowbreak{}a\allowbreak{}m\allowbreak{}p\allowbreak{}l\allowbreak{}e\allowbreak{}s\allowbreak{}/\allowbreak{}d\allowbreak{}o\allowbreak{}c\allowbreak{}s}} is now a living product surface, not just a showcase shell
\item
  the next DOGFOOD work is quality and depth, not pretending there is a hidden extra family-count milestone
\end{itemize}

This is important for both users and maintainers. Users get a much more credible field guide. Maintainers get a stronger proving surface for shell, design-language, and docs changes.

\section{Localization: Real Packages And Workflows}\label{record-275-localization-real-packages-and-workflows}

Localization moved from strategy-doc ambition into actual packages and workflows.

The relevant landed work includes:

\begin{itemize}
\tightlist
\item
  \href{https://github.com/flyingrobots/bijou/blob/cbd70612140e7a450d764b3c7be1c6af878b99bc/docs/design/LX-001-bijou-i18n-runtime-package.md}{LX-001}
\item
  \href{https://github.com/flyingrobots/bijou/blob/cbd70612140e7a450d764b3c7be1c6af878b99bc/docs/design/LX-002-bijou-i18n-tools-package.md}{LX-002}
\item
  \href{https://github.com/flyingrobots/bijou/blob/cbd70612140e7a450d764b3c7be1c6af878b99bc/docs/design/LX-003-spreadsheet-adapters-and-catalog-exchange-workflows.md}{LX-003}
\item
  \href{https://github.com/flyingrobots/bijou/blob/cbd70612140e7a450d764b3c7be1c6af878b99bc/docs/design/LX-004-provider-adapters-for-workbook-and-bundle-exchange.md}{LX-004}
\item
  \href{https://github.com/flyingrobots/bijou/blob/cbd70612140e7a450d764b3c7be1c6af878b99bc/docs/design/LX-005-rich-spreadsheet-and-filesystem-adapters.md}{LX-005}
\item
  \href{https://github.com/flyingrobots/bijou/blob/cbd70612140e7a450d764b3c7be1c6af878b99bc/docs/design/LX-006-xlsx-localization-adapters.md}{LX-006}
\item
  \href{https://github.com/flyingrobots/bijou/blob/cbd70612140e7a450d764b3c7be1c6af878b99bc/docs/design/LX-008-localize-shell-chrome-and-dogfood.md}{LX-008}
\end{itemize}

The workspace now includes:

\begin{itemize}
\tightlist
\item
  \texttt{@\allowbreak{}f\allowbreak{}l\allowbreak{}y\allowbreak{}i\allowbreak{}n\allowbreak{}g\allowbreak{}r\allowbreak{}o\allowbreak{}b\allowbreak{}o\allowbreak{}t\allowbreak{}s\allowbreak{}/\allowbreak{}b\allowbreak{}i\allowbreak{}j\allowbreak{}o\allowbreak{}u\allowbreak{}-\allowbreak{}i\allowbreak{}1\allowbreak{}8\allowbreak{}n}
\item
  \texttt{@\allowbreak{}f\allowbreak{}l\allowbreak{}y\allowbreak{}i\allowbreak{}n\allowbreak{}g\allowbreak{}r\allowbreak{}o\allowbreak{}b\allowbreak{}o\allowbreak{}t\allowbreak{}s\allowbreak{}/\allowbreak{}b\allowbreak{}i\allowbreak{}j\allowbreak{}o\allowbreak{}u\allowbreak{}-\allowbreak{}i\allowbreak{}1\allowbreak{}8\allowbreak{}n\allowbreak{}-\allowbreak{}t\allowbreak{}o\allowbreak{}o\allowbreak{}l\allowbreak{}s}
\item
  \texttt{@\allowbreak{}f\allowbreak{}l\allowbreak{}y\allowbreak{}i\allowbreak{}n\allowbreak{}g\allowbreak{}r\allowbreak{}o\allowbreak{}b\allowbreak{}o\allowbreak{}t\allowbreak{}s\allowbreak{}/\allowbreak{}b\allowbreak{}i\allowbreak{}j\allowbreak{}o\allowbreak{}u\allowbreak{}-\allowbreak{}i\allowbreak{}1\allowbreak{}8\allowbreak{}n\allowbreak{}-\allowbreak{}t\allowbreak{}o\allowbreak{}o\allowbreak{}l\allowbreak{}s\allowbreak{}-\allowbreak{}n\allowbreak{}o\allowbreak{}d\allowbreak{}e}
\item
  \texttt{@\allowbreak{}f\allowbreak{}l\allowbreak{}y\allowbreak{}i\allowbreak{}n\allowbreak{}g\allowbreak{}r\allowbreak{}o\allowbreak{}b\allowbreak{}o\allowbreak{}t\allowbreak{}s\allowbreak{}/\allowbreak{}b\allowbreak{}i\allowbreak{}j\allowbreak{}o\allowbreak{}u\allowbreak{}-\allowbreak{}i\allowbreak{}1\allowbreak{}8\allowbreak{}n\allowbreak{}-\allowbreak{}t\allowbreak{}o\allowbreak{}o\allowbreak{}l\allowbreak{}s\allowbreak{}-\allowbreak{}x\allowbreak{}l\allowbreak{}s\allowbreak{}x}
\end{itemize}

And the shell/docs surface now has real localized behavior instead of only English-first assumptions.

Important release note: the i18n packages are now part of the planned \texttt{v\allowbreak{}4\allowbreak{}.\allowbreak{}1\allowbreak{}.\allowbreak{}0} automated release path. The publish workflow, release dry-run workflow, and release docs now all agree that these packages are public release artifacts for the \texttt{4\allowbreak{}.\allowbreak{}1\allowbreak{}.\allowbreak{}0} line.

\section{Design Language: More Explicit Product Taste}\label{record-275-design-language-more-explicit-product-taste}

Bijou's design language is much less implicit than it was at \texttt{v\allowbreak{}4\allowbreak{}.\allowbreak{}0\allowbreak{}.\allowbreak{}0}.

The post-\texttt{v\allowbreak{}4\allowbreak{}.\allowbreak{}0\allowbreak{}.\allowbreak{}0} work established:

\begin{itemize}
\tightlist
\item
  explicit UX doctrine and content guidance
\item
  a calmer explainability surface
\item
  inspector rhythm and inspector drawer patterns
\item
  stronger shell review-surface patterns
\item
  more concrete design-system guidance across component families
\end{itemize}

Representative cycles include:

\begin{itemize}
\tightlist
\item
  \href{https://github.com/flyingrobots/bijou/blob/cbd70612140e7a450d764b3c7be1c6af878b99bc/docs/design/DL-001-capture-design-language-doctrine.md}{DL-001}
\item
  \href{https://github.com/flyingrobots/bijou/blob/cbd70612140e7a450d764b3c7be1c6af878b99bc/docs/design/DL-003-prove-canonical-patterns-in-shared-surfaces.md}{DL-003}
\item
  \href{https://github.com/flyingrobots/bijou/blob/cbd70612140e7a450d764b3c7be1c6af878b99bc/docs/design/DL-004-prove-drawer-rhythm-and-notice-rows.md}{DL-004}
\item
  \href{https://github.com/flyingrobots/bijou/blob/cbd70612140e7a450d764b3c7be1c6af878b99bc/docs/design/DL-005-prove-inspector-and-guided-flow-rhythm.md}{DL-005}
\item
  \href{https://github.com/flyingrobots/bijou/blob/cbd70612140e7a450d764b3c7be1c6af878b99bc/docs/design/DL-006-prove-inspector-panel-rhythm.md}{DL-006}
\item
  \href{https://github.com/flyingrobots/bijou/blob/cbd70612140e7a450d764b3c7be1c6af878b99bc/docs/design/DL-007-promote-inspector-panel-block.md}{DL-007}
\end{itemize}

The result is that Bijou now has stronger answers to questions like:

\begin{itemize}
\tightlist
\item
  what should shell review surfaces feel like?
\item
  how should AI reasoning be surfaced honestly?
\item
  what does calm terminal hierarchy look like?
\item
  which UI choices are doctrine, and which are merely possible?
\end{itemize}

\section{Performance, Validation, And Release Hardening}\label{record-275-performance-validation-and-release-hardening}

The repo also got materially stronger operationally.

Notable improvements since \texttt{v\allowbreak{}4\allowbreak{}.\allowbreak{}0\allowbreak{}.\allowbreak{}0}:

\begin{itemize}
\tightlist
\item
  renderer benchmark coverage and low-allocation renderer work
\item
  stronger smoke tooling and example/workflow validation
\item
  a DOGFOOD coverage gate
\item
  a design-system docs preflight
\item
  a workflow shell preflight
\item
  a repo-native release readiness script
\item
  improved publish and release-dry-run workflow compatibility
\end{itemize}

This work is less glamorous than the runtime and shell work, but it matters because it makes the release more reproducible and the repo less dependent on luck.

\section{Repo Signposts And Engineering Doctrine}\label{record-275-repo-signposts-and-engineering-doctrine}

The repository itself is easier to navigate and reason about than it was at \texttt{v\allowbreak{}4\allowbreak{}.\allowbreak{}0\allowbreak{}.\allowbreak{}0}.

New or newly formalized repo surfaces include:

\begin{itemize}
\tightlist
\item
  \href{https://github.com/flyingrobots/bijou/blob/5b10cb89165b76673ffc60da59d7de5a531e55e9/docs/METHOD.md}{METHOD}
\item
  \href{https://github.com/flyingrobots/bijou/blob/5b10cb89165b76673ffc60da59d7de5a531e55e9/docs/BEARING.md}{BEARING}
\item
  \href{https://github.com/flyingrobots/bijou/blob/cbd70612140e7a450d764b3c7be1c6af878b99bc/docs/VISION.md}{VISION}
\item
  \href{https://github.com/flyingrobots/bijou/blob/cbd70612140e7a450d764b3c7be1c6af878b99bc/docs/system-style-javascript.md}{System-Style JavaScript}
\item
  \href{https://github.com/flyingrobots/bijou/blob/cbd70612140e7a450d764b3c7be1c6af878b99bc/docs/release.md}{Release Guide}
\end{itemize}

This is mostly maintainer-facing, but it has a user-visible effect: the repo is less likely to drift between code, docs, backlog, and release claims.

\section{What This Release Is Best At}\label{record-275-what-this-release-is-best-at}

Compared to \texttt{v\allowbreak{}4\allowbreak{}.\allowbreak{}0\allowbreak{}.\allowbreak{}0}, this release is better when you need:

\begin{itemize}
\tightlist
\item
  a clearer fullscreen application substrate
\item
  a calmer framed shell with more truthful interaction behavior
\item
  a serious living docs surface instead of a demo stub
\item
  first-class localization packages and workflow seams
\item
  clearer guidance about how Bijou wants to be used
\end{itemize}

\section{What This Release Does Not Pretend}\label{record-275-what-this-release-does-not-pretend}

This release still does not mean:

\begin{itemize}
\tightlist
\item
  the runtime-engine migration is completely finished
\item
  every shell concern has been pushed fully behind the runtime seams
\item
  DOGFOOD is ``done forever''
\item
  localization is complete across every shell surface
\end{itemize}

The point is not that the work is finished. The point is that the repo now has much more honest runtime truth, shell truth, docs truth, and release truth than it had at \texttt{v\allowbreak{}4\allowbreak{}.\allowbreak{}0\allowbreak{}.\allowbreak{}0}.

